\section{Control, sensing, and programmability}
\label{sec:control}

Programmability is the distinguishing experimental resource of
robotic active matter. In conventional active systems, propulsion
and interactions are often fixed by particle chemistry, geometry,
or mechanical properties. Robotic systems permit us to change
local rules, sensing, communication, feedback, and sometimes
the interactions themselves. Hence, a collective state can be
investigated not only by observing what emerges, but also by
changing the mechanism believed to produce it.

Yet, programmability introduces a fundamental ambiguity. A
controller may realize the microscopic assumptions of a theory,
but it may also impose the collective behavior that the theory
is intended to explain. A flock maintained by explicit
alignment feedback is not equivalent to polar order generated
by mechanical collisions. Similarly, a rotating assembly
driven by a common magnetic field is not necessarily a
spontaneously rotating active phase.

We therefore consider control as both an experimental advantage
and a potential confounder. Its scientific value depends on
which physical variables become independently adjustable and
which collective observables remain genuine predictions.

The extraction distinguishes five control architectures:
no control or passive rules, on-board rules, external feedback,
central control, and learned policies. These categories describe
the primary architecture assigned to each study, not every
control mechanism present in an experiment. In particular,
on-board control can coexist with communication, while
external feedback may involve a central measurement system
without directly specifying every agent's trajectory.

The October 17, 2026 evidence base contains
491 included studies. Control-architecture
counts are calculated from the corrected
extraction table. The categories describe
the primary control architecture assigned
to each study and are interpreted alongside
the corresponding platform and interaction
characteristics.

Figure~\ref{fig:control} relates control architecture to
the maximum number of agents used in hardware experiments.
The diagram highlights an important distinction between
the complexity of a controller and the scale of the
physical collective. Neither determines the scientific
quality of an experiment by itself.

\begin{figure}[t]
\centering
\includegraphics[width=\columnwidth]{figures/subtaxonomy_control}
\caption{Control architectures and experimental collective
sizes in robotic active matter. The bars distinguish
no control or passive rules, on-board rules, external
feedback, central control, and learned policies.
Colors indicate the largest reported hardware
collective in each study, grouped into 2--9,
10--99, 100--999, and at least 1000 agents.
The size groups are descriptive, not measures of
physical scalability or statistical independence.
Studies with missing agent maxima and control
classifications require separate consideration.
The graphic derives from the original extraction
snapshot, before one duplicate microscopic
record was merged.}
\label{fig:control}
\end{figure}

\subsection{No control or passive rules}
\label{sec:control-passive}

The absence of an explicit feedback controller does not
mean that a robotic collective lacks interactions. Indeed,
mechanically driven systems may be especially informative
precisely because their collective response is not
programmed at the level of neighbor decisions. Individual
activity is generated by an actuator, while collective
organization arises through contact, friction, fluid
coupling, geometry, and confinement.

The extraction assigns 72 studies to no control or
passive rules. Vibration-driven robots form the largest
group, with 41 studies. Other implementations contribute
13, microscopic robotic agents eleven, active solids
and modular robots two each, and smarticles and
wheeled or legged systems the remaining three.

This classification should not be interpreted as
absence of external energy input. An actively driven
robot is not passive in the thermodynamic sense.
Rather, the label indicates that the extracted
primary architecture does not use explicit feedback
to prescribe collective behavior.

\paragraph*{What becomes experimentally accessible?}

Mechanically driven robots offer direct access to
the relationship between individual propulsion and
interaction-mediated motion. Their collective states
include clustering, transport, boundary accumulation,
rotation, and spatial organization
\cite{Welch2018aggregatesuppressio,Welch2019activereorientation,Deblais2019howtoforceanarmy,Xi2024beamsbotsandbey,Balda2022playingwithactive,Xi2024emergentintelligenc,Mohanty2025multimodalmotionan,GuzmnLastra2026emergenceofbeating,Edwards2026robusttopologically,Musacchio2026activeembracemente,Escobar2025effectofparticlec,Scholz2018rotatingrobotsmove,Arbel2025amechanicalroutef,Antonov2025selfsustainedfrict,Baconnier2022selectiveandcollec,Boudet2025nonmarkoviandynami}.

The important advantage is that collective order need not be
encoded in a controller. If collisions produce alignment
or crowding produces spatial organization, those effects
must arise through the physical implementation. This
provides a useful contrast to systems that explicitly
calculate desired neighbor-dependent velocities.

Mechanical activity can nevertheless involve complex
single-agent dynamics. Vibration may produce correlated
rotational fluctuations, intermittent contact, and
direction-dependent motility. These effects can determine
the observed collective state. The strongest experiments
therefore characterize isolated-agent propulsion and
pairwise interactions before interpreting the collective.

Additional mechanically driven studies investigate
confinement, transport, fluctuations, and collective
rearrangements
\cite{MuozObreque2026dynamicsofmagnetic,Xi2024emergentbehaviorso,Tian2018flowratebehavioran,Porvatov2021optimizingselfrota,Giomi2013swarmingswirlinga,Luo2025flockingphasesepar,Musacchio2026fluidizationinduced,Deblais2018boundariescontrolc,Boudet2022effectivetemperatur,Engbring2023anonlinearfluctuat,Dmitriev2021statisticalcorrelat,Callegari2023playingwithactive,Callegari2024activematterexperi,Balda2024playingwithactive,Li2021programmingactivec,Siebers2023exploitingcompositi,Son2024emergentfunctional,Boudet2021fromcollectionsof,Wagner2026geometryinducedtra,Lvay2026collectivedynamics}.
These studies are relevant to active Brownian and
other stochastic descriptions because the physical
actuation can be compared with measured collective
transport or spatial correlations.

At microscopic scales, control-free or passively
coupled dynamics occur under different physical
conditions. The extraction includes field-mediated
and mechanically interacting systems
\cite{Calabr2026interactiondrivena,Rivas2023spatialpatterningo,Cheng2025selfassemblyandco,Yang2022emergentmicroroboti,Urso2023reconfigurableself,Yan2024programmingstructur,Yigit2020cohesiveselforgani,Desgranges2023microswimmersunder,Ceron2023programmableselfor,Wang2018collectivesofspinn,Zhang20233dprintedmicromete}.
The common point is not an identical propulsion
mechanism. It is that the extracted collective
behavior need not arise from an explicit
state-dependent control policy.

\paragraph*{Where interpretation becomes difficult.}

The absence of feedback removes one possible
source of artificially imposed collective order.
It does not make the interaction law transparent.

In the 72 studies, steric contact is recorded
in 66, environment-mediated interactions in
39, magnetic coupling in 18, and hydrodynamic
coupling in twelve. These modalities often
occur together. Thus, an observed collective
state may be compatible with several
mechanistic explanations.

For example, boundary accumulation may depend
on persistent propulsion, wall collisions,
substrate friction, or particle geometry.
A rotating cluster may reflect intrinsic
chirality, magnetic forcing, or
interaction-induced collective circulation.
Without independently varying or measuring
these ingredients, the experiment cannot
uniquely identify its mechanism.

The relevant studies include mechanically
and environmentally coupled systems
\cite{Zhou2021synchronizedswimmin,Zion2024amechanicalorigin,Bo2026threephasesofodd,Hu2026phototunablenonreci,Yang2025electronicfreepart,Newbolt2024flowinteractionsle,Vantomme2021coupledliquidcryst,Boudet2025effectsofboundarie,Das2025roleofkinematicco,Warmington2026adiscontinuouslyco,B2023adiscontinuouslyco,Xia2024biomimeticsynchroni,Miyashita2007selfassembleofwat,Mayya2019nonuniformrobotde,Zhang2024proteinfoldinginsp,Miyashita2008howmorphologyaffec,Ma2026crosslinkcollectiv,Li2013asymmetricgearrect,Leech2026spincoupledhydrody}.
They demonstrate the range of physical mechanisms
available without a conventional distributed
feedback architecture.

\paragraph*{Evidence quality.}

Of the 72 studies, 34 report quantitative
theory--experiment agreement, 30 qualitative
agreement, seven no comparison, and one has
a missing agreement value.

For D6, 22 studies are rated low risk,
31 have some concerns, and 19 are high risk.
This distribution is comparatively favorable
for experiments in which interactions often
emerge through physical coupling.

The replication profile is less strong:
14 studies have low-risk D1 ratings,
53 have some concerns, and five are high risk.
Independent-run counts are recorded in
16 studies. Explicit code availability
occurs in seven, and data availability
in eleven. Ten studies were assessed
on abstracts alone.

These findings suggest that passive-rule
and mechanically driven systems provide
valuable quantitative tests, but incomplete
replication information limits confidence
in several reported effects.

Their main contribution is the ability
to examine collective behavior without
explicitly programming a desired
macroscopic state. Their main limitation
is reduced independence between
propulsion, interactions, and
environmental conditions.

Hence, the best experimental strategy
is to characterize the physical
interaction law and vary confinement,
density, or driving systematically.
This would make passive-rule systems
stronger tests of active-matter theory
rather than merely simpler
alternatives to programmable swarms.

\subsection{On-board rules}
\label{sec:control-onboard}

On-board rules are the dominant control architecture
in the extracted evidence. In these systems,
individual agents determine their behavior from
locally available information. The information
may come from sensors, neighboring robots,
internal state, or a combination of these.

The extraction contains 278 studies with on-board
rules. Wheeled or legged swarm robots contribute
168, followed by 61 residual platforms,
19 vibration-driven systems, eleven modular
robots, nine smarticles, seven active solids,
two microscopic robotic systems, and one
record without a platform classification.

The prevalence of on-board control reflects a
fundamental strength of robotics: microscopic
interaction rules can be implemented in
individually operating physical agents.
This makes such systems especially suitable
for experimental realizations of distributed
collective-motion models.

\paragraph*{Programmable local interactions.}

A conventional distributed controller computes
an agent's next action from its available state
and information about its surroundings. In a
general notation, this can be written as
\begin{equation}
\mathbf u_i(t)=
\pi_i\!\left[
\mathbf x_i(t),
\mathcal I_i(t);
\boldsymbol{\lambda}
\right],
\label{eq:control-onboard-policy}
\end{equation}
where $\mathbf x_i$ is the local state,
$\mathcal I_i$ denotes accessible information,
and $\boldsymbol{\lambda}$ contains
adjustable parameters.

The physical significance of this expression
lies in the distinction between the controller
output and the resulting collective motion.
The policy determines how an agent attempts to
move. Its realized trajectory also depends
on actuation, collisions, sensing errors,
and interactions with the environment.

The included studies investigate alignment,
collective motion, aggregation, and programmed
coordination through local rules
\cite{Miki2010asimplerulebased,Chen2013aflockingalgorithm,Wang2021fielddriverobotsw,Fruchart2022nonreciprocalphase,Yang2023collectivechaseem,Paramanick2024emergenceofrunand,Kernbach2011threecasesofconne,Tran2022frontierledswarmin,Giardina2022collectivephototact,Vega2023simulatelessexpec,Vega2023indirectswarmcontr,Vega2024agentbasedemulatio,Varadharajan2026programmableflocks,Gao2023robustdistributedc,Lee2022distributedcohesive,Garnier2005aggregationbehaviou,Baumann2024hybridflockforma,Guzzi2014bioinspiredobstacle,Turgut2008modelingphasetrans,Navarro2008anapproachtoflock,Ercan2011swarmrobotflocking,Ferrante2013aselfadaptivecomm,Turgut2008selforganizedflocka,Vsrhelyi2018optimizedflockingo}.

In these systems, programmability can separate
the intended interaction rule from the
observed physical collective response.
For example, one may change the effective
alignment range without rebuilding the
robots. One may also introduce directional
interactions, delays, or changes in
neighbor selection.

Such interventions are useful because they
can discriminate between theoretical
mechanisms. An alignment model predicts
not only motion in a common direction,
but also the dependence of collective
order on interactions and fluctuations
\cite{Vicsek1995novel}
\cite{Toner1995longrange}.

The difficulty is that the rule itself
can make the outcome unsurprising.
If every robot calculates a target
velocity aligned with its neighbors,
observing alignment confirms that the
hardware broadly realizes the controller.
A stronger physical test concerns the
resulting fluctuations, stability,
correlations, and transitions.

Studies investigating collective response,
group deformation, indirect interaction,
and changes in local information are
particularly relevant
\cite{Li2016turinglearningam,Lee2022distributeddeformab,Karagzel2022collectivegradient,Zhang2022collectivefissionb,Xiao2023perceptionofmotion,Zheng2024bodyorientationcha,Zheng2026nestedinteractionn,Zhou2021collectivesynchroni,Mijalkov2016engineeringsensoria,Noirhomme2024brainbotsassmarta,Paramanick2025spontaneousemergenc,Barone2024experimentalandnum,Zhou2025activesearchpromot,Liu2026noiseenabledgoala,Dizenhaus2026anonymousleadership,Wang2025hyperuniformnetwork}.
Their extracted observables permit different
degrees of confrontation with physical theory.
However, the presence of a measured collective
quantity does not establish that a theoretical
prediction was tested without fitting.

\paragraph*{Distributed control beyond mobile swarms.}

On-board rules are not restricted to wheeled robots.
Modular and granular systems can use internal
actuation to generate collective configurations
or locomotion
\cite{Li2018supersmarticlealo,Bhamla2021collectivetransport,Vardhan2022glidersinshapecha,Ma2024smarticle20desig,Yang2024selfamputatingand,Hu2026diatominspired1di,Zhang2026anisotropicshapesh,Li2019particleroboticsba,Oliveri2021continuouslearning,Drew2021acousticcommunicati,Savoie2019arobotmadeofrobo,Chvykov2020lowrattlingapred,Li2020programmingactiveg}.

Active structures provide another example.
Programmed local actuation can modify
mechanical coupling and collective response
\cite{Brandenbourger2019nonreciprocalrobot,Brandenbourger2021limitcyclesturnac,Binysh2026moreislessinunpe,Du2026metamaterialsthatl,Coulais2025activemetamaterials,Chen2020selfsensingmetamat}.
The relevant physical variables are then
deformation, mechanical response,
and structural dynamics rather than
flocking polarization alone.

This reveals an important generalization.
On-board rules need not represent
social or informational interactions
between freely moving agents.
They can also determine actuation
within a mechanically connected
structure. The resulting theory
should describe the appropriate
physical degrees of freedom.

\paragraph*{Scale and robustness.}

Among the 278 on-board-rule studies,
153 have a recorded maximum of two
to nine agents, 93 reach ten to 99,
eight reach 100 to 999, and six
reach at least 1000. Eighteen
have no recorded agent maximum.

Thus, this architecture supports
both small experiments and large
collectives. Yet, most reported
systems are modest in size.
A large number of independently
controlled robots is not by itself
evidence for large-system
active-matter scaling.

The main challenge is to test
whether collective observables
change systematically with
system size, rather than merely
demonstrating that a controller
continues to function.

For theoretical comparison,
50 studies report quantitative
agreement, 145 qualitative agreement,
and 83 no comparison. D6 is rated
low risk in 34 studies,
some concerns in 72,
and high risk in 172.

Hence, despite the prevalence
of programmable local rules,
strong quantitative theory
comparisons remain a minority.

Replication reporting is
more favorable but incomplete.
Eighty-seven studies receive
low-risk D1 ratings,
183 have some concerns,
and eight are high risk.
Independent-run counts
are recorded for 133 studies.
Explicit code availability
occurs in 52, and data
availability in 28.
Fourteen studies were
assessed from abstracts alone.

These findings are consistent
with a literature that has
developed substantial capability
for implementing distributed
collective behavior, but
has not yet converted
that capability systematically
into predictive experiments
on active matter.

The next step is to exploit
the independent control
of interaction parameters.
A well-characterized
on-board implementation
could vary sensing range,
delay, alignment strength,
or interaction symmetry
while measuring several
collective observables.

This would distinguish
robust physical consequences
from controller-specific
behavior. It would also
provide more stringent
tests than demonstrating
that robots successfully
execute a collective task.

\subsection{External feedback}
\label{sec:control-external}

External feedback modifies agent behavior using
measurements obtained outside the individual agent.
A camera or other measurement system may estimate
positions, orientations, or collective state, and
the resulting information may determine subsequent
actuation.

This architecture is particularly useful when
the physical agents are too small to carry
sensors or processors. It also permits
interaction rules that would be difficult
to realize through direct mechanical forces.

The extraction assigns 84 studies to external
feedback. Microscopic robotic agents contribute
49 and feedback-controlled colloids twelve.
Wheeled or legged robots contribute twelve,
other platforms four, modular robots three,
vibration-driven robots three,
and active solids one.

The distribution shows that external feedback
is not merely another way to coordinate
macroscopic robot swarms. It is a central
experimental mechanism for making microscopic
active systems programmable.

\paragraph*{Feedback as an effective interaction.}

A feedback-controlled particle may change its
motion according to the measured configuration
of neighboring particles. Schematically,
we can express this as
\begin{equation}
\mathbf u_i(t)=
\mathcal F_i\!\left[
\{\mathbf r_j(t-\tau)\},
\boldsymbol{\lambda}
\right],
\label{eq:control-feedback}
\end{equation}
where $\tau$ represents measurement and
actuation latency.

This formulation makes an important physical
point. Feedback can generate an effective
interaction even when no corresponding
conservative force exists between the
particles. Such interactions may depend
on orientation, configuration, or
collective state.

The included colloidal studies investigate
programmable rotation, configuration-dependent
response, and collective organization
\cite{Martinet2022rotationcontrolin,Yan2014colloidalsuperstruc,Chen2024synchronousandfull,Buerle2020formationofstable,Aubret2021metamachinesofplur,Heuthe2025tunablecolloidalsw,Knippenberg2022roleofcohesionin,Han2020programmablechiral,Chen2022collectiveresponse,Guo2026afeedbackcontrolle,Chen2024persistentandrespo,Lffler2021behaviordependentc}.

Their experimental strength is the ability
to change the effective interaction law
while retaining comparable particles.
This permits tests of collective states
under controlled changes in coupling.

Yet, the feedback loop becomes part of
the active-matter dynamics. Measurement
noise, update frequency, and latency
can influence the resulting motion.
A theory of instantaneous interactions
may therefore fail even when the
nominal feedback rule resembles it.

The strongest experiments characterize
the implemented effective interaction
independently and then compare
collective observables with its
predictions. This is especially
important when the feedback
algorithm and the theoretical
model share the same mathematical
structure.

\paragraph*{Microscopic control and collective response.}

Microscopic robot experiments use external
feedback and programmable fields to
produce collective transport,
self-assembly, reconfiguration,
and rotational motion
\cite{Sun2021swarmingmicrodrople,Xie2020bioinspiredsoftmic,MayorgaMartinez2023multimodaldrivenma,Yigit2019programmablecollect,Das2023optoelectronictraje,Feng2023weakionexchangeba,Cao2024mofbasedmagneticm,Li2023swarmingresponsive,Gong2025collectivereconfigu,Chen2020temporallightmodul,Chen2022programmablespatio,Xu2022collectivebehaviors,Liu2023swarmingmultifuncti,Chen2023nanorobotmediateds,Ussia2024magneticmicrorobot,Tran2025roboticmicrocapsule,Xia2025artificialintellige,Wang2026reconfigurableassem,Kim2026magneticmxenebased,Aubret2018targetedassemblyan,Ahmed2021bioinspiredacousto,Raghavendra2026programmablepersist}.
These experiments extend control toward
physical regimes dominated by magnetic,
hydrodynamic, and other microscopic
interactions.

An important distinction is whether
the applied field is uniform or
spatially structured. A common
field may synchronize propulsion
without generating an independently
programmable interaction between agents.
Feedback based on individual or
collective configuration offers
a more direct route to
state-dependent coupling.

Other included microscopic studies
investigate controlled assembly,
field-mediated transport,
and collective reconfiguration
\cite{Gardi2022microrobotcollectiv,Alapan2019shapeencodeddynami,Wang2018collectivebehavior,Ji2021ondemandassemblya,An2024modelfreecontrolo,Fan2023largescaleswarmco,Xie2019programmablegenerat,Wang2022reconfigurablemagne,Wang2022magneticactuationo,Law2024navigationofmagnet,Yang2019statisticsbasedaut,Yu2022adaptivepatternand,Yu2019reconfigurablecollo,Wang2021precisecontrolofm,Ceron2026fluidictorqueenabl,Dong2020controllingtwodime,Dong2019collectiveformation,Li2022vectorcontrolledwh}.
Their collective observables can establish
transport and spatial organization,
but a physical interpretation requires
knowledge of how the external input
affects individual motion and
interparticle coupling.

\paragraph*{Evidence quality.}

In the 84 external-feedback studies,
34 report quantitative agreement,
31 qualitative agreement, and
19 no theory comparison.

For D6, 21 studies are rated
low risk, 25 have some concerns,
and 38 are high risk.
For replication, fifteen
receive low-risk D1 ratings,
67 have some concerns,
and two are high risk.

Independent-run counts are
recorded in 29 studies.
Explicit code availability
occurs in five and data
availability in eight.
Three were assessed
from abstracts alone.

Agent maxima are missing
for 36 studies. Among the
remaining records,
21 have maxima of two
to nine agents,
14 reach ten to 99,
ten reach 100 to 999,
and three reach at least
1000 agents.

These values require care
because a microscopic
particle population and
a collective of individually
addressable macroscopic
robots are not directly
comparable.

The evidence supports
external feedback as a
powerful method for
implementing programmable
collective interactions.
However, incomplete
reporting of feedback
implementations and
experimental repetitions
limits independent
validation.

The most important
next step is therefore
not simply finer
control of trajectories.
It is calibration
of the entire
measurement--feedback
loop, followed by
quantitative tests
of collective response
under independently
varied control
parameters.



\subsection{Central control}
\label{sec:control-central}

Central control makes collective behavior accessible through
a shared decision process. A controller may observe the state
of an entire robotic assembly and determine the actuation
of individual agents or groups. This architecture differs
from distributed on-board rules because interactions need
not be computed locally. Instead, the controller can use
global information to construct the effective dynamics.

For active-matter physics, this is a powerful experimental
capability. One can impose spatially heterogeneous fields,
modify collective response, or introduce interactions that
would be difficult to realize through direct mechanical
coupling. Yet, the same capability creates a fundamental
interpretive problem. If the controller continuously
specifies the collective configuration, the observed
organization may be imposed rather than emergent.

The extraction assigns 43 studies to central control.
Wheeled or legged robots contribute 16, microscopic
robotic agents 14, modular or self-reconfigurable
robots seven, and other platforms six. All 43
records were assessed on a full-text basis.

The principal reported phenomenon is self-assembly
or shape formation, present in 30 studies.
Transport or locomotion occurs in 22,
flocking or polar order in eleven,
chirality or rotation in nine,
and clustering and synchronization
in eight each. These labels overlap.

The distribution shows that central control is
particularly relevant to configuration and
transport. Its main contribution is not
necessarily to establish spontaneous ordering.
Rather, it makes collective geometry and
dynamical response experimentally programmable.

\paragraph*{What does global information add?}

Consider a collective state vector
$\mathbf X(t)$ containing the measured
positions, orientations, or other
degrees of freedom. A central controller
may prescribe actions according to
\begin{equation}
\mathbf U(t)=
\Pi\!\left[
\mathbf X(t);
\boldsymbol{\lambda}
\right].
\label{eq:control-central-policy}
\end{equation}
Here, $\Pi$ is the collective control
law and $\boldsymbol{\lambda}$ denotes
adjustable parameters.

The scientific advantage is that
$\Pi$ can implement interactions
that depend on the global state.
For example, one can vary the
strength or geometry of effective
coupling while holding the
physical hardware approximately
constant.

This permits controlled studies
of collective response. However,
the resulting system is not
equivalent to one with purely
local interactions. The central
controller may introduce
information transfer over
distances that exceed the
physical interaction range.

Hence, the range and symmetry
of the effective interaction
must be described as properties
of the experimental model,
not hidden implementation
details.

Macroscopic robotic experiments
investigate collective motion,
configuration, and controlled
coordination through such
architectures
\cite{Yamada2011circularmotiongene,Sankar2026trendformationwith,Pan2026exploringthephysic,Zhou2024anticipationpromote,Bahaidarah2024swarmflockingusing,Santos2014cohesionandsegrega,PliegoJimnez2023flockingandformati,Hou2009ontheimplementatio,Min2011designandanalysis,Beattie2025realizingemergentc,Li2026onfeedbackspeedco,Zheng2024emergenceofcollect,Zhou2025gazingpreferencein,Zheng2025shapingcollectiveb,Maffettone2024mixedrealityenviro,Abpeikar2023iterativetransferl}.
These implementations support experimental
investigations of coordinated movement,
collective structure, and responses
under programmable interactions.

Nevertheless, the physical conclusion
depends on the observables that are
measured. A stable target formation
establishes successful collective
control. A phase diagram,
correlation function, or
measured response to perturbations
can provide stronger evidence
about the resulting active dynamics.

\paragraph*{Global actuation at microscopic scales.}

Central control is particularly useful
for microscopic robotic agents.
Individual particles generally cannot
carry conventional computing and
communication hardware. Instead,
external actuation can determine
their motion through magnetic,
optical, acoustic, or other
physical mechanisms.

The included studies examine controlled
assembly, collective motion, and
reconfiguration
\cite{Pauer2021programmabledesign,McGovern2024multifunctionalchir,Sun2024multiplemagnetoopt,Je2024spatiallyselective,Li2023swarmingresponsive,Liu2022bionicmorphological,Chen2020temporallightmodul,Chen2022programmablespatio,Tran2025roboticmicrocapsule,Baker2019shapeprogrammed3d,Gao2025softmagneticmicror,Raghavendra2026programmablepersist,Ahmed2021bioinspiredacousto,Chen2024magneticactuatedhy}.

These experiments can make the
driving conditions highly
reproducible in principle.
The field geometry and temporal
protocol can be specified
independently of the
observed collective state.

Yet, a spatially varying field
may directly generate the
observed collective pattern.
We must therefore distinguish
a field-imposed configuration
from interaction-mediated
organization.

A useful test compares
the predicted response
of individual particles
with the collective
response under the
same driving conditions.
Deviations can reveal
the contribution of
interparticle interactions.

Furthermore, external
field control may
modify propulsion
and coupling simultaneously.
In magnetic assemblies,
for example, changes
in the driving conditions
may alter individual
movement and relative
particle interactions.
A quantitative model
must account for both.

\paragraph*{Connected and modular systems.}

Central control also appears
in modular or reconfigurable
assemblies
\cite{Yi2025reconfigurablerobot,Zhao2024snailinspiredrobot,Bhattacharjee2026collectiveassembly,Kaeser2025individualandcolle,Ding2025transformablemodula,Zhao2022softlatticemodules,Paulos2015automatedselfassem}.
These experiments investigate
configuration and collective
locomotion through coordinated
components.

Here, the main interpretive
distinction concerns control
and mechanical emergence.
A prescribed sequence
of module movements
can generate a complex
collective deformation.
However, the same outcome
may depend on contact
forces and connectivity
that are not explicitly
specified by the controller.

This creates an opportunity
to study how imposed
actuation interacts
with physical constraints.
The most informative
measurements include
relative module motion,
deformation, transport,
and response under
changes in connectivity.

Other central-control
implementations address
collective formation,
flocking, and transport
\cite{PitaRomero2026flockingbehaviorfo,Jia2014experimentalimpleme,Fu2023tabletopplatformo,Delight2016developingrobotics,Su2026safetycriticalmult,MendozaRobles2026distributedformatio}.
They further illustrate
that centralized decision-making
can be implemented through
different physical platforms.

\paragraph*{Evidence quality and limitations.}

Among the 43 centrally controlled
studies, eight report quantitative
theory--experiment agreement,
28 qualitative agreement,
and seven no comparison.

For D6, six studies receive
low-risk ratings, 15 have
some concerns, and 22
are high risk.
Thus, most studies
do not provide
low-risk quantitative
tests of a theoretical
collective response.

Replication information is
also incomplete.
Nine studies are rated
low risk for D1,
33 have some concerns,
and one is high risk.
Independent-run counts
are recorded for twelve
studies.

Explicit code availability
is identified in five studies,
and data availability in five.
Maximum agent counts
are available for 35 studies.
Twenty report maxima
of two to nine agents,
and fifteen reach
ten to 99.
The remaining eight
have no recorded maximum.

These numbers do not
indicate that central
control is restricted
to small physical
systems in general.
They describe the
included experiments
and the completeness
of their extraction.

The evidence supports
central control as
an effective means
of producing and
modifying collective
behavior. Yet, its
capacity to impose
the desired state
requires particular
care in theoretical
interpretation.

We therefore distinguish
controlled realization
from explanatory validation.
A central controller
can implement a proposed
interaction law accurately.
The physical test concerns
whether the resulting
collective observables
follow predictions that
were not directly
specified by that law.

Future experiments
should report the
actual control inputs,
their dependence on
collective state,
and the resulting
physical response.
A particularly useful
design would compare
centralized and
distributed implementations
of equivalent
effective interactions.

Such comparisons
could establish
which collective
features depend
on information
architecture and
which depend on
the resulting
physical coupling.

\subsection{Learned policies}
\label{sec:control-learning}

Learning introduces another route to programmable
collective behavior. Rather than specifying every
interaction rule analytically, a policy may be
obtained through optimization, imitation, or
reinforcement. Its actions can depend on local
observations, collective state, or an internally
learned representation.

For robotics, this provides a flexible means
of generating behavior. For active-matter physics,
the implications are more subtle. A learned policy
may discover collective responses not anticipated
by a hand-designed controller. Yet, the resulting
interaction law may be difficult to interpret.

The extraction identifies eleven studies with
a learned policy as their primary control
architecture. Six involve wheeled or legged
robots, three other platforms, one a robotic
micro- or nano-agent system, and one
feedback-controlled active colloids.

Seven studies report flocking or polar order,
five swarming or vortex behavior,
four transport or locomotion,
three chirality or rotation,
and two clustering.
The phenomena overlap within studies.

The small size of this category is important.
Learning is represented in the experimental
evidence, but it is not yet a dominant
control architecture among the included
physical collectives.

\paragraph*{Learning collective interaction rules.}

The relevant experiments investigate
learned flocking, collective motion,
milling, aggregation, and
environment-dependent coordination
\cite{Schilling2019learningvisionbase,Choi2026communicationfreec,Kraus2026generativeadversari,Wu2024adversarialimitatio,Lffler2023collectiveforaging,Liu2025collectivemotionmo,An2026ondemandpatterntr,Zhang2025learningefficientf,Zhu2025amillingswarmofg,Zhu2025fromspikestoswarm,Jia2025adeepreinforcement}.

These studies demonstrate several ways
in which learning can contribute
to physical collective behavior.
A policy may map sensory information
to motion, adapt interactions,
or generate a collective response
under changing conditions.

The scientific advantage is that
the effective interaction need
not be prescribed in advance.
If a learned controller
produces an unexpected
collective pattern, the
experiment may suggest
new mechanisms for
theoretical investigation.

However, the policy itself
does not explain the
resulting behavior.
An observed flock or
vortex may arise through
a complicated mapping
from observations to actions.
Without further analysis,
the physical interaction
responsible for the
collective state remains
underdetermined.

This is a particularly
important distinction
between learning for
robotic performance
and learning as an
experimental tool
for active-matter physics.

A policy that achieves
a collective objective
may be useful even when
its mechanism is opaque.
A scientific explanation
requires additional
measurements linking
observations, actions,
and collective response.

\paragraph*{Programmability beyond hand-designed rules.}

Learning can explore interaction
strategies that would be
difficult to enumerate
manually. This may
include nonlinear responses,
state-dependent couplings,
or context-sensitive
changes in behavior.

Yet, a learned policy
may also exploit
features specific
to the training
environment.
Such behavior
can produce excellent
experimental performance
without supporting
a general physical
principle.

For example, a policy
trained in one arena
may rely on boundary
cues. A controller
optimized for one
population size
may behave differently
when the number
of agents changes.
A response learned
under one sensing
configuration may
not transfer to
another.

These possibilities
are scientifically
testable. They
suggest experiments
in which the
training conditions
and physical
evaluation conditions
are separated.

A particularly
informative study
would evaluate
a learned interaction
under changes
in density,
boundary geometry,
sensor noise,
or collective size
that were absent
from training.

If the resulting
collective behavior
remains consistent
with a simple
physical model,
the experiment
may reveal a
generalizable
effective interaction.

If it does not,
the failure can
identify which
environmental
or informational
details determine
the learned state.

\paragraph*{Evidence quality.}

Of the eleven learned-policy
studies, one reports
quantitative theoretical
agreement, four qualitative
agreement, and six
no comparison.

None receives a
low-risk D6 rating.
Three have some concerns,
and eight are high risk.

For replication, one
study receives a
low-risk D1 rating,
while ten have
some concerns.
Independent-run counts
are recorded in six
studies.

Explicit code availability
occurs in two studies,
and data availability
in two.
All eleven were
assessed on full text.

Maximum agent counts
are recorded for ten
studies. Five have
maxima of two to nine,
and five reach
ten to 99 agents.
One maximum is missing.

These findings indicate
that learned policies
have demonstrated
physical collective
behavior, but the
evidence for quantitative
active-matter theory
remains limited.

In particular,
the presence of
a learned collective
state does not
establish that
the policy has
discovered a
specific physical
interaction law.

The most productive
research direction
is therefore to
combine learning
with interpretable
physical measurements.

The learned policy
should be analyzed
through controlled
interventions:
which observations
change its actions,
which interactions
are effectively
implemented, and
which collective
quantities remain
stable under
distribution shifts?

Such experiments
could connect
data-driven control
to active-matter
theory without
assuming that
successful learned
behavior is itself
a physical
explanation.

\subsection{Sensing, communication, and information-mediated interactions}
\label{sec:control-information}

Control architecture specifies where decisions are made.
Interaction modality specifies how agents influence
one another. These dimensions overlap, but they
should not be confused. An on-board controller may
use communication, local sensing, or mechanical
contact. A central controller may reproduce an
interaction that would otherwise be implemented
locally.

For active-matter physics, this distinction
matters because information itself can mediate
effective interactions. Agents may change
their motion in response to signals,
observations, or communicated states,
even when no corresponding mechanical
force acts between them.

Such systems are often described by
effective alignment, attraction,
repulsion, or response laws.
Yet, the information-processing
mechanism can introduce asymmetry,
delay, finite perception range,
and state dependence that are
absent from simple particle models.

\paragraph*{Local sensing as an interaction mechanism.}

The extracted interaction vocabulary identifies
150 studies with optical or other sensing.
These are predominantly on-board systems:
121 use on-board rules, fourteen external
feedback, nine learned policies, five
central control, and one passive rules.

The category includes experiments on
visual alignment, collective coordination,
sensing-dependent motion, and
perception-mediated interactions
\cite{Schilling2019learningvisionbase,Choi2026communicationfreec,Wu2024adversarialimitatio,Liu2025collectivemotionmo,Zhu2025amillingswarmofg,Zhu2025fromspikestoswarm,Jia2025adeepreinforcement,Mijalkov2016engineeringsensoria,Castro2025visualcollectivebe,Zhou2025gazingpreferencein,Zheng2024emergenceofcollect,Zheng2025shapingcollectiveb,Li2026onfeedbackspeedco,Zhang2025learningefficientf,Kraus2026generativeadversari}.

The physical effect of sensing
depends on what is measured.
A robot may estimate
neighbor position,
relative velocity,
orientation, or
the structure of
its environment.

These quantities
can generate different
effective interaction
laws even when the
resulting collective
motion looks similar.

For example,
orientation-dependent
perception can produce
asymmetric responses.
One agent may detect
another without
being detected
in return.
The resulting
coupling is then
nonreciprocal
at the level
of the controller.

This creates
a useful experimental
route to investigating
nonreciprocal collective
dynamics. However,
the implementation
must be characterized
carefully.

Camera geometry,
occlusion,
measurement frequency,
and finite processing
time may affect
the effective interaction.
If these features
are omitted from
the theoretical model,
agreement may depend
on accidental
compensation between
parameters.

The most informative
experiments therefore
measure the response
of an individual
robot to controlled
sensory inputs before
examining collective
behavior.

This permits the
effective interaction
law to be estimated
independently.
Collective observables
can then provide
a genuine test
of the resulting
model.

\paragraph*{Communication and collective coupling.}

Communication is recorded as an interaction
modality in 156 studies. Of these, 132 use
on-board rules, sixteen central control,
six external feedback, and two
learned policies.

This distribution reflects the importance
of information exchange in conventional
swarm-robotics experiments.

The included literature contains
collective coordination,
communication-dependent flocking,
distributed interaction rules,
and changes in collective
configuration
\cite{Cianci2007communicationinas,Christensen2008swarmorphscripta,Ferrante2013aselfadaptivecomm,Valentini2015collectivedecision,Kernbach2011threecasesofconne,Drew2021acousticcommunicati,Cazenille2024hearingtheshapeof,Gao2023robustdistributedc,Lee2022distributedcohesive,Zhou2024anticipationpromote,Turgut2008selforganizedflocka,Vsrhelyi2018optimizedflockingo,Tran2022frontierledswarmin,Maffettone2024mixedrealityenviro,Santos2014cohesionandsegrega,Min2011designandanalysis,Beattie2025realizingemergentc,MendozaRobles2026distributedformatio}.

Communication permits agents
to respond to information
that is not available
through direct physical
contact.
It can therefore
increase the effective
interaction range
or alter the
topology of coupling.

This makes communication
a physical control
parameter in a
broader sense.
Changing which
agents exchange
information can
change the collective
state without
changing mechanical
hardware.

However, a communicated
state is not equivalent
to a mechanical
interaction force.
The resulting
collective dynamics
depend on how
the received information
is converted into
actuation.

Communication delay
is another important
difference. In
an instantaneous
interaction model,
agents respond
to the current
state of their
neighbors.
A communicating
robot may respond
to an earlier
measurement.

This delay can
alter stability,
oscillations,
and collective
response.
Hence, the
effective theory
should include
delay when it
is dynamically
relevant.

Communication
also raises
the question
of network topology.
A group can
interact through
a fixed information
graph even while
its physical
positions change.
Conversely,
distance-dependent
communication may
generate a graph
that evolves
with collective
motion.

These alternatives
need not produce
the same active
dynamics.
A meaningful
comparison should
therefore distinguish
spatial proximity
from information
connectivity.

\paragraph*{Programmed rules and physical interactions.}

The interaction vocabulary identifies
161 studies containing programmed rules.
Among them, 120 use on-board rules,
nineteen external feedback,
seventeen central control,
four passive rules,
and one a learned policy.

These counts demonstrate that
programmed interactions are not
restricted to one control
architecture. Nevertheless,
the extracted label does not
specify whether the rule
implements alignment,
attraction, repulsion,
state-dependent switching,
or another interaction.

The relevant physical interpretation
therefore requires examination
of the actual controller
and its observables.

The theoretical opportunity
is considerable.
A programmed interaction
can be changed systematically
while keeping the same
physical platform.

For example, one could
compare reciprocal and
nonreciprocal alignment
or vary the effective
interaction range.
One could also introduce
controlled delay
or modify the response
to local density.

These interventions
would permit experiments
that distinguish
competing theoretical
predictions.

Yet, a change in
programmed rules
may also alter
the collective
operating regime
through actuator
constraints or
sensing limitations.
The nominal software
parameter is therefore
not always identical
to the effective
physical parameter.

The relevant experiment
should measure
the implemented
response law
rather than assume
perfect realization
of the controller.

\paragraph*{Quality of information-mediated evidence.}

The extracted agreement ratings
show that 28 of the
150 sensing-labelled
studies report
quantitative theory
comparison, while
70 report qualitative
comparison and 52
no comparison.

For communication,
25 of 156 studies
report quantitative
comparison, 86
qualitative comparison,
and 45 no comparison.

Programmed-rule studies
contain 36 quantitative,
84 qualitative,
and 41 without
a recorded comparison.

These modality
categories overlap.
Their counts must
not be added
to calculate
the number of
distinct studies.

Risk-of-bias
ratings reveal
the same limitation.
Among sensing-labelled
studies, eighteen
receive low-risk
D6 ratings,
25 have some concerns,
and 107 are high risk.

For communication,
the corresponding
counts are seventeen,
43, and 96.
For programmed rules,
they are 22,
47, and 92.

Thus, the widespread
use of information-mediated
interactions has not
yet produced an
equally widespread
body of low-risk
quantitative
theory tests.

The pattern is
consistent with
the broader
control-architecture
analysis. Robotic
systems offer
many ways to
program collective
behavior, but
implementation
success is often
more thoroughly
demonstrated than
the predictive
validity of the
corresponding
physical model.

\subsection{Synthesis: control as an experimental variable}
\label{sec:control-synthesis}

The comparison of control architectures establishes
a central opportunity for robotic active matter.
We can intervene directly on microscopic
interaction rules while observing collective
physical consequences.

Yet, this opportunity is realized differently
across experimental platforms.

Mechanically driven systems provide collective
behavior without explicitly programmed
neighbor responses. Their advantage is that
interactions arise from physical coupling.
Their limitation is the difficulty of
independently varying propulsion,
noise, friction, and contact dynamics.

On-board rules provide direct implementation
of local interactions. Their advantage is
programmability at the agent level.
Their limitation is that the controller
may impose the behavior whose emergence
the experiment is intended to explain.

External feedback permits effective
interactions in systems too small
for conventional on-board control.
Its advantage is flexible coupling
based on measured collective state.
Its limitation is that sensing,
latency, and feedback become
part of the physical dynamics.

Central control makes collective
configuration and global interaction
rules accessible. Its advantage
is precise intervention.
Its limitation is the possibility
of replacing spontaneous
organization with externally
prescribed coordination.

Learned policies extend
programmability beyond
hand-designed interaction
rules. Their advantage
is the exploration
of new collective
responses.
Their limitation
is the difficulty
of identifying
the effective
physical mechanism.

These distinctions
show that control
architecture is
not simply an
engineering property.
It can determine
which theoretical
description is
appropriate.

\paragraph*{The controller belongs to the physical model.}

A common source of interpretive error
is to treat control as an implementation
detail while attributing collective
behavior entirely to physical interactions.

This separation is often invalid.
An agent that changes direction
based on sensed neighbors
experiences an effective
interaction created
by its controller.
A feedback-controlled
particle can exhibit
configuration-dependent
forces that do not
exist without
the measurement
and actuation loop.

Consequently,
a faithful physical
model must include
the relevant features
of the control
architecture.

This does not mean
that every controller
detail must appear
in a continuum
theory.
Coarse-graining
is precisely the
process of identifying
which microscopic
details remain
important at
collective scales.

However,
that simplification
must be tested,
not assumed.

For example,
if communication
delay has no
measurable effect
on collective
order within
a specified regime,
an instantaneous
interaction model
may be adequate.
If delay changes
stability or
fluctuations,
it must be
included in
the effective
description.

The same reasoning
applies to
sensing anisotropy,
actuator limitations,
and common
external forcing.

\paragraph*{Programmability as causal intervention.}

The strongest use of control
is to vary a proposed
mechanism while measuring
its collective consequences.

For example,
an experiment
might hold
agent propulsion
approximately
constant while
changing the
interaction range.
Alternatively,
it might preserve
the effective
alignment rule
while changing
the sensing
modality.

Such comparisons
would help determine
whether a collective
state depends
on information
architecture
or on the
effective
interaction law.

A further
possibility
is to construct
different microscopic
implementations
of the same
predicted collective
dynamics.

If a mechanical
system, a communicating
robot swarm,
and a feedback-controlled
colloidal assembly
exhibit comparable
response under
appropriately matched
parameters, the
result would provide
stronger evidence
for generality
than repeated
experiments
using one
controller.

Yet, matching
platforms requires
careful choice
of observables
and dimensionless
control parameters.
A similar
trajectory pattern
alone does not
establish equivalent
physics.

\paragraph*{What should be reported?}

The evidence suggests a practical
standard for future experiments.

An on-board or centralized
robotic study should
document the implemented
interaction law,
including relevant
sensing and
communication constraints.

A feedback-controlled
experiment should
characterize the
measurement process,
control latency,
and mapping from
observations
to actuation.

A learned-policy
experiment should
separate training
from evaluation
and report tests
under changed
physical conditions.

A mechanically driven
experiment should
measure isolated-agent
motion and
interaction-dependent
response.

Across all architectures,
collective observables
should be reported
with independent
repetitions and
appropriate uncertainty.
Code and data
availability would
further support
verification and
comparative analysis.

The evidence does
not imply that
one architecture
is universally
superior.
Rather, each
offers a different
combination of
experimental
control and
physical interpretability.

The most decisive
experiments will
exploit these
complementary
strengths.

\paragraph*{From control performance to physical understanding.}

A central conclusion of this
section is that programmable
collective behavior and
predictive active-matter
physics are distinct
scientific achievements.

Robotic experiments
have established
many ways to
generate flocking,
aggregation,
collective transport,
synchronization,
and shape formation.

Yet, the extracted
quality assessments
show that low-risk
quantitative
theory comparisons
remain less common
than successful
collective
implementations.

The next advance
should therefore
focus on experiments
in which control
parameters serve
as independent
physical interventions.

We should ask
not only whether
a controller
produces the
desired behavior,
but how the
collective state
changes when
the interaction
is modified,
whether the
response agrees
with theory,
and whether
that agreement
survives changes
in platform
and scale.

This perspective
also gives learning
a productive
scientific role.
A learned policy
may generate
an unexpected
collective response.
We can then
analyze its
effective interaction,
construct a
simpler physical
model, and
test the model
through controlled
perturbations.

Thus, learning
and programmability
need not compete
with physical
interpretability.
They can provide
new hypotheses
and new interventions,
provided that
the resulting
mechanisms are
subsequently tested.

We conclude that
control is most
valuable to
active-matter
physics when
it converts
a theoretical
assumption
into an
experimentally
adjustable
quantity.

The robot
then becomes
more than
a physical
demonstration
of collective
behavior.
It becomes
an instrument
for testing
which microscopic
interactions
are necessary,
which are
sufficient,
and which
remain relevant
at collective
scales.

